% ----------------------------------------------------------------------------
% Resumen / Abstract
% ----------------------------------------------------------------------------

\noindent
El cáncer de mama es el tumor de mayor incidencia entre las mujeres a nivel
mundial, y la mamografía sigue siendo la prueba de cribado de referencia para su
detección precoz. En los últimos años, las técnicas de \textit{deep learning}
han demostrado un gran potencial en el análisis automático de imagen médica,
pero suelen comportarse como cajas negras y son sensibles al ruido y a la
variabilidad de adquisición. El \textbf{Análisis Topológico de Datos} (TDA, por
sus siglas en inglés) ofrece una perspectiva complementaria: cuantifica la
\emph{forma} de los datos mediante invariantes topológicos robustos frente al
ruido, como la homología persistente.

\medskip
\noindent
Este Trabajo de Fin de Grado estudia la aplicación del TDA a la clasificación y
visualización de mamografías digitales del conjunto de datos público
\mbox{CBIS-DDSM}. Se extraen descriptores topológicos (diagramas de
persistencia, códigos de barras e imágenes de persistencia) a partir de las
imágenes y de las regiones de interés, se evalúa su capacidad discriminativa
entre lesiones benignas y malignas, y se comparan y combinan con modelos de
aprendizaje profundo. El objetivo último es obtener un sistema de apoyo al
diagnóstico más \emph{interpretable} sin sacrificar rendimiento.

\vspace{1cm}
\noindent\textbf{Palabras clave:} análisis topológico de datos, homología
persistente, mamografía, cáncer de mama, aprendizaje profundo, imagen médica,
CBIS-DDSM.

\bigskip\bigskip
\begin{center}\rule{0.5\textwidth}{0.4pt}\end{center}
\bigskip

\selectlanguage{english}
\noindent\textbf{Abstract.}
Breast cancer is the most common malignancy among women worldwide, and
mammography remains the reference screening test for its early detection. Deep
learning has shown remarkable potential in medical image analysis, yet such
models often behave as black boxes and are sensitive to noise and acquisition
variability. \textbf{Topological Data Analysis} (TDA) provides a complementary
viewpoint by quantifying the \emph{shape} of the data through noise-robust
topological invariants such as persistent homology. This bachelor's thesis
investigates the application of TDA to the classification and visualisation of
digital mammograms from the public CBIS-DDSM dataset. Topological descriptors
(persistence diagrams, barcodes and persistence images) are extracted and
evaluated for their ability to discriminate benign from malignant lesions, and
compared and combined with deep learning models, aiming at a more
\emph{interpretable} decision-support system.

\medskip
\noindent\textbf{Keywords:} topological data analysis, persistent homology,
mammography, breast cancer, deep learning, medical imaging, CBIS-DDSM.
\selectlanguage{spanish}
